\subsection{Project Structure}
\label{sec:impl_structure}

The codebase follows a feature-first layout in which every domain is a self-contained package under \texttt{abridgeai/features/}. Each feature package contains the same internal layers: \texttt{models.py} (SQLAlchemy ORM), \texttt{schemas/} (Pydantic DTOs), \texttt{queries/} (database access functions), \texttt{services/} (business logic), \texttt{routers/} (FastAPI route handlers), and optionally \texttt{workers/} (background tasks) and \texttt{ai/} (AI pipeline stages). This uniform layout makes it straightforward to locate any piece of logic and to add new features without touching unrelated code.

\begin{figure}[H]
\centering
\begin{verbatim}
abridgeai/
+-- api/                  # Application entry point, router registration
+-- core/                 # Cross-cutting: config, db, cache, audit, security
|   +-- audit/            # Async audit-log listener and worker actor
|   +-- cache/            # Redis client, decorators, cache-key registry
|   +-- db/               # Soft-delete mixin, conflict mapper, hard-delete guard
|   +-- observability/    # Structured logging, HTTP audit log
|   +-- pagination/       # Cursor-based pagination
|   +-- security/         # JWT creation, token hashing, security headers
+-- ai/                   # Shared AI primitives
|   +-- chunking/         # Token-aware, semantic, timestamp-aware chunkers
|   +-- extraction/       # Multimodal extractors (PDF, DOCX, audio, video, ...)
|   +-- knowledge_graph/  # Neo4j builder and retrieval
|   +-- llm/              # LLM gateway, embedding client, pricing, audit
|   +-- prompts/          # Jinja2 prompt loader
|   +-- retrieval/        # pgvector, BM25, RRF fusion, MMR, role filter
+-- features/
|   +-- access_control/   # Organisations, org units, RBAC
|   +-- admin/            # Admin dashboards, AI cost tracking, audit viewer
|   +-- career_paths/     # Career path authoring and learner enrollment
|   +-- courses/          # Course, module, lesson authoring and learner view
|   +-- enrollments/      # Enrollment lifecycle management
|   +-- identity/         # Users, auth sessions, Google OAuth, MFA
|   +-- interviews/       # AI interview generation, session taking, evaluation
|   +-- materials/        # Learning material upload, ingestion pipeline
|   +-- notifications/    # In-app and email notification dispatch
|   +-- progress/         # Lesson progress tracking and analytics
|   +-- quizzes/          # AI quiz generation, taking, grading
|   +-- spaced_repetition/# SM-2 scheduler, card review, dashboards
+-- infrastructure/       # Third-party adapters: Neo4j, S3/Garage, Google OAuth
\end{verbatim}
\caption{Top-level directory layout of the aBridgeAI backend.}
\label{fig:project_structure}
\end{figure}

\paragraph{Layered dependency rule}
Dependencies flow strictly downward: routers depend on services, services depend on queries and AI primitives, queries depend on models. No layer imports from a layer above it. Cross-feature communication is mediated through the \texttt{api/public} module of each feature, which exposes a stable internal API surface. This prevents circular imports and keeps feature boundaries explicit.

\paragraph{Database migrations}
Schema evolution is managed by 16 Alembic migration files numbered \texttt{0001} through \texttt{0016}. Migration \texttt{0001} establishes the full baseline schema ported from a hand-managed SQL file. Subsequent migrations add partial unique indexes, cascade rules, audit columns, pgvector half-precision embeddings (\texttt{halfvec(3072)}), and full-text search vectors. Every migration is reversible via a \texttt{downgrade()} function.
